%==============================================================================
%  CV — Oscar Ordoñez
%
%  Réplica del formato original (Word/Calibri) medida directamente del PDF:
%    · Papel Letter, márgenes 0.75 in izq. / 0.72 in der.
%    · Calibri 20 / 11.5 / 10.6 / 10 / 9.5 / 9 pt
%    · Azul oscuro #1F3864, azul medio #2E5496, gris #595959
%    · Regla de 0.7 pt bajo cada encabezado de sección
%
%  Cambios frente al original: los enlaces ahora son hipervínculos reales
%  (LinkedIn, GitHub, portafolio, OdobGames, correo y teléfono) y se añadió
%  el sitio nuevo https://odobgames.github.io/Portafolio/
%
%  La etapa en BBVA aparecía como un solo cargo de Septiembre 2024 a Enero
%  2026. Los primeros seis meses fueron la práctica profesional en el área de
%  Rentabilidad, así que ahora son dos entradas separadas:
%    · Practicante · Rentabilidad          Septiembre 2024 -- Febrero 2025
%    · Data Analyst Specialist · BI        Marzo 2025      -- Enero 2026
%
%  Se quitó el \newpage manual que forzaba el salto de página: con una entrada
%  más el corte ya no caía donde debía. Revisa dónde parte ahora la página y,
%  si quieres fijarlo, vuelve a poner \newpage donde te convenga.
%
%  Compilar con pdfLaTeX (Overleaf: Menu -> Compiler -> pdfLaTeX) o con XeLaTeX/Tectonic.
%==============================================================================
\documentclass[letterpaper,10pt]{article}

% Carlito is metrically compatible with Calibri: same widths, same look, and unlike
% Calibri it is available on Overleaf. The "lining" option forces lining figures: with
% oldstyle figures the digits are not mapped to Unicode and an ATS would read the
% phone, the dates and the email as garbage.
\usepackage{iftex}
\ifPDFTeX
  \usepackage[T1]{fontenc}
  \usepackage[utf8]{inputenc}
\fi
\usepackage[sfdefault,lining]{carlito}

\usepackage[spanish,es-noshorthands]{babel}
\usepackage[letterpaper,top=0.83in,bottom=0.65in,left=0.75in,right=0.72in]{geometry}
\usepackage{xcolor}
\usepackage{enumitem}
\usepackage{needspace}
\usepackage[hidelinks]{hyperref}

% Glyph-to-Unicode map: with pdfLaTeX it makes the text that an ATS (or a copy and
% paste) reads come out with accents and without odd ligatures. Other engines skip it.
\ifdefined\pdfgentounicode
  \input{glyphtounicode}
  \pdfgentounicode=1
\fi

%------------------------------------------------------------------ colores --
\definecolor{cvdark}{RGB}{31,56,100}   % #1F3864  nombre y cargos
\definecolor{cvblue}{RGB}{46,84,150}   % #2E5496  subtítulo, secciones, reglas
\definecolor{cvgray}{RGB}{89,89,89}    % #595959  contacto, fechas, empresa

% Los enlaces conservan el gris del original para no alterar el aspecto.
% ¿Los quieres visibles en azul? Cambia cvgray por cvblue en urlcolor.
\hypersetup{
  colorlinks=true, urlcolor=cvgray, linkcolor=cvgray,
  pdftitle={CV — Oscar Ordoñez},
  pdfauthor={Oscar Ordoñez},
  pdfkeywords={Gobierno de Datos, Databricks, Unity Catalog, AWS, Python, SQL, PySpark, DAMA-DMBOK2, LLM, Calidad de Datos},
  pdflang={es-CO},
  pdfsubject={Ingeniero de Gobierno de Datos · Científico de Datos · Ingeniero de IA}
}

%--------------------------------------------------------------- tipografía --
\newcommand{\sizeName}   {\fontsize{20}{23}\selectfont}
\newcommand{\sizeTitle}  {\fontsize{10.6}{12.7}\selectfont}
\newcommand{\sizeContact}{\fontsize{9}{11}\selectfont}
\newcommand{\sizeSection}{\fontsize{11.5}{13.8}\selectfont}
\newcommand{\sizeBody}   {\fontsize{10}{12.2}\selectfont}
\newcommand{\sizeWhere}  {\fontsize{9.5}{11.4}\selectfont}
\newcommand{\sizeSkill}  {\fontsize{9.5}{11.6}\selectfont}

\newcommand{\md}{\,\textperiodcentered\,}                  % separador ·
\newcommand{\sepdot}{\hspace{4pt}\textbullet\hspace{4pt}}   % separador •

\setlength{\parindent}{0pt}
\pagestyle{empty}

% Word no divide palabras con guion: justifica estirando los espacios. Estas
% líneas reproducen ese comportamiento y evitan cortes como "Data Wa-".
\hyphenpenalty=10000
\exhyphenpenalty=10000
\tolerance=9999
\emergencystretch=3em

%--------------------------------------------------------------- estructura --
% Encabezado de sección: versalitas azules + regla de 0.7 pt a todo lo ancho.
\newcommand{\cvsection}[1]{%
  \needspace{3\baselineskip}%
  \vspace{13pt}%
  {\color{cvblue}\sizeSection\bfseries #1}\par
  \vspace{-5pt}%
  {\color{cvblue}\rule{\textwidth}{0.7pt}}\par
  \vspace{3pt}%
}

% Cargo a la izquierda, fechas alineadas a la derecha, en la misma línea.
\newcommand{\cvrole}[2]{%
  \needspace{3\baselineskip}%
  {\color{cvdark}\sizeTitle\bfseries #1}%
  \hfill{\color{cvgray}\sizeContact #2}\par
  \vspace{1pt}%
}

% Empresa y ubicación, en cursiva gris.
\newcommand{\cvwhere}[1]{%
  {\color{cvgray}\sizeWhere\itshape #1}\par
  \vspace{3pt}%
}

% Viñetas: marca a 8 pt del margen, texto a 18 pt (medido del original).
\setlist[itemize]{
  leftmargin=18pt, labelwidth=10pt, labelsep=0pt, align=left,
  topsep=1pt, itemsep=3pt, parsep=0pt, partopsep=0pt,
  label=\textbullet, before=\raggedright
}

% Línea de habilidades: categoría en negrita azul, contenido en negro.
\newcommand{\cvskill}[2]{%
  {\raggedright\sizeSkill{\color{cvdark}\bfseries #1:} #2\par}%
  \vspace{3pt}%
}

%==============================================================================
\begin{document}
\sizeBody

%--------------------------------------------------------------- encabezado --
{\centering
  {\color{cvdark}\sizeName\bfseries OSCAR ORDOÑEZ}\\[3pt]
  {\color{cvblue}\sizeTitle Ingeniero de Gobierno de Datos\md Científico de Datos\md Ingeniero de IA}\\[1pt]
  {\color{cvgray}\sizeContact
    Bogotá, Colombia\sepdot
    \href{tel:+573227596832}{+57 322 759 6832}\sepdot
    \href{mailto:13odob27@gmail.com}{13odob27@gmail.com}}\\[0pt]
  {\color{cvgray}\sizeContact
    \href{https://www.linkedin.com/in/oscarordonez13}{LinkedIn}\sepdot
    \href{https://github.com/OdobGames}{GitHub}\sepdot
    \href{https://odobgames.github.io/Portafolio/}{Portafolio}\sepdot
    \href{https://odobgames.github.io/OdobGames}{OdobGames}}\par}
\vspace{0pt}

%---------------------------------------------------------------- perfil -----
\cvsection{PERFIL PROFESIONAL}

Científico de la Computación de la Universidad Nacional de Colombia,
especializado en Gobierno de Datos, ingeniería de datos en la nube (Databricks,
AWS) e Inteligencia Artificial aplicada, con experiencia en los sectores Telco
(Claro Colombia) y FinTech (BBVA). Enfocado en el gobierno del ciclo de vida del
dato y la auditoría de uso de activos de información en ecosistemas de gran
escala, alineado con los marcos DAMA-DMBOK2, DCAM v2.2 e ISO 25012. Experto en
Python, SQL, PySpark, en plataformas Databricks (Unity Catalog, Delta Lake) y
Teradata, y en la integración de LLMs para la automatización de procesos de
datos. Inglés avanzado (C1) y experiencia bajo metodologías ágiles (Scrum).

%------------------------------------------------------------ experiencia ---
\cvsection{EXPERIENCIA PROFESIONAL}

\cvrole{Ingeniero de Gobierno de Datos}{Abril 2026 -- Presente}
\cvwhere{Grupo CINTE --- en misión en Claro Colombia\md Bogotá}
\begin{itemize}
  \item \textbf{Gobierno del ciclo de vida del dato:} Implementación de prácticas
        de Data Governance sobre un ecosistema de Data Warehouse multiplataforma,
        asegurando la trazabilidad, calidad y estandarización de los activos de
        información a lo largo de su ciclo de vida.
  \item \textbf{Auditoría de uso de activos:} Diseño y desarrollo de soluciones
        analíticas en Python para la auditoría de uso del inventario de datos y
        la identificación de objetos redundantes, como insumo para las políticas
        corporativas de depuración y optimización.
  \item \textbf{Catalogación con IA:} Aplicación de LLMs y técnicas de muestreo
        estadístico para automatizar la catalogación, clasificación y
        enriquecimiento semántico de metadatos a gran escala.
  \item \textbf{Políticas y estándares de gobierno:} Definición de políticas de
        uso, acceso e incidentes para entornos analíticos en Databricks (Unity
        Catalog), alineadas con DAMA-DMBOK2, DCAM v2.2 e ISO 25012, bajo
        metodología Scrum.
\end{itemize}
\vspace{7pt}

\cvrole{Desarrollador de Inteligencia Artificial (Contratista)}{Enero 2026 -- Febrero 2026}
\cvwhere{Softgic\md Colombia}
\begin{itemize}
  \item \textbf{Desarrollo de agente de IA:} Construcción de un agente de
        Inteligencia Artificial end-to-end, integrando LLMs con flujos de
        automatización para resolver casos de uso de negocio.
  \item \textbf{Workflows en n8n:} Diseño, codificación y orquestación de
        workflows complejos en n8n, integrando APIs, fuentes de datos y servicios
        de IA para automatizar procesos críticos.
  \item \textbf{SQL y modelado de datos:} Desarrollo intensivo en SQL para
        extracción, transformación y análisis de datos como soporte al
        comportamiento del agente.
  \item \textbf{Metodologías ágiles:} Entregas iterativas bajo marcos Scrum, con
        ciclos cortos de validación y despliegue.
\end{itemize}
\vspace{7pt}

\cvrole{Data Analyst Specialist\md Business Intelligence --- Client Solutions}{Marzo 2025 -- Enero 2026}
\cvwhere{BBVA Colombia\md Bogotá}
\begin{itemize}
  \item \textbf{Arquitectura de modelo de rentabilidad (P\&L):} Diseño e
        implementación del modelo de rentabilidad de la unidad adquirente sobre
        el ecosistema AWS (SageMaker, Athena, S3), procesando grandes volúmenes
        de datos transaccionales.
  \item \textbf{Pipelines ETL/ELT:} Desarrollo de pipelines de datos automatizados
        con PySpark y SQL avanzado para optimizar la ingesta, transformación y
        disponibilidad de información para las áreas de negocio.
  \item \textbf{Business Intelligence y visualización:} Diseño de dashboards
        estratégicos en MicroStrategy y Google Sheets (Apps Script) para el
        seguimiento de los KPIs clave del negocio adquirente.
  \item \textbf{Modelado predictivo:} Construcción de modelos predictivos con
        Python (Pandas, Scikit-learn) para soporte a la toma de decisiones
        financieras y comerciales.
  \item \textbf{IA Generativa:} Investigación e integración de LLMs y APIs de IA
        Generativa para apoyar el análisis avanzado y la generación automatizada
        de reportes ejecutivos.
\end{itemize}
\vspace{7pt}

\cvrole{Practicante\md Rentabilidad --- Client Solutions}{Septiembre 2024 -- Febrero 2025}
\cvwhere{BBVA Colombia\md Bogotá}
\begin{itemize}
  \item \textbf{Automatización de procesos operativos:} Desarrollo de rutinas de
        automatización para las tareas operativas recurrentes del área de
        Rentabilidad, sustituyendo procedimientos que hasta entonces se
        ejecutaban de forma manual.
  \item \textbf{Soporte a la operación de banca minorista:} Apoyo a la operación
        diaria del negocio de banca minorista dentro de la unidad de Client
        Solutions.
\end{itemize}
\vspace{7pt}

\cvrole{Desarrollador de Software, IA \& Videojuegos (Indie)}{Enero 2023 -- Presente}
\cvwhere{OdobGames\md Proyecto personal}
\begin{itemize}
  \item \textbf{Desarrollo orientado a datos:} Aplicación de IA y analítica para
        mejorar productos digitales y optimizar la experiencia de usuario.
  \item \textbf{Full-Stack Mobile \& Game Development:} Desarrollo móvil con Unity
        (C\#) y prototipado en Unreal Engine (UEFN), gestionando el ciclo de vida
        completo: programación, build y publicación en tiendas.
  \item \textbf{Automatización \& APIs:} Implementación de flujos en n8n e
        integración de APIs para la optimización de procesos internos del estudio.
\end{itemize}

%-------------------------------------------------------------- educación ---
\cvsection{EDUCACIÓN}

\cvrole{Ciencias de la Computación (B.S.)}{2019 -- 2025}
\cvwhere{Universidad Nacional de Colombia\md Bogotá, Colombia}

%------------------------------------------------------------- habilidades --
\cvsection{HABILIDADES TÉCNICAS}

\cvskill{Gobierno de Datos}{Data Governance, Data Catalog, Data Quality, Metadata
  Management, Data Lineage, Master Data Management, DAMA-DMBOK2, DCAM v2.2,
  ISO 25012.}

\cvskill{Cloud \& Big Data}{Databricks (Unity Catalog, Delta Lake, Spark), AWS
  (SageMaker, Athena, S3), Teradata, PySpark, Apache Spark, Data Warehouse.}

\cvskill{Ciencia de Datos \& IA}{Python (Pandas, NumPy, Scikit-learn), Modelado
  Predictivo, Machine Learning, LLMs, GenAI, Prompt Engineering, Integración de
  APIs de IA.}

\cvskill{Bases de Datos \& ETL}{SQL avanzado, Modelado de Datos, ETL/ELT,
  Optimización de Queries, Data Profiling, Procesamiento Distribuido.}

\cvskill{Business Intelligence}{MicroStrategy, Google Sheets (Apps Script),
  Visualización de Datos, KPIs, Dashboards Ejecutivos.}

\cvskill{Automatización \& Agentes IA}{n8n, Workflows Orchestration, AI Agents,
  Integración de APIs.}

\cvskill{Desarrollo de Software}{Python, C\# (Unity), Unreal Engine (UEFN), Git,
  GitHub, Azure DevOps.}

\cvskill{Metodologías}{Scrum, Agile, Trabajo Colaborativo en Equipos
  Multidisciplinarios.}

\cvskill{Idiomas}{Español (Nativo) \textbar{} Inglés Avanzado (C1).}

\end{document}
